\documentclass[11pt]{article}
\usepackage[margin=1in]{geometry}
\usepackage{times}
\usepackage[T1]{fontenc}
\usepackage{lmodern}
\usepackage[hidelinks]{hyperref}
\setlength{\parskip}{6pt}
\setlength{\parindent}{0pt}

\begin{document}

\begin{center}
{\large \textbf{Aligned but Not Retrieved: Target-Language Selection Failure\\
in Bilingual Contrastive Fine-Tuning}}\\[4pt]
isaacchang@g.harvard.edu, ingridchien@g.harvard.edu, mtran@mit.edu
\end{center}

\vspace{4pt}

\textbf{Motivation}

Multilingual retrieval-augmented generation (RAG) systems face a failure mode we call
language drift: a user queries in English targeting a Spanish knowledge base, but the
retriever---a sentence encoder fine-tuned with contrastive learning---returns French or
German documents instead, because semantically equivalent translations are
indistinguishable in embedding space, and the language model generates in the wrong
language. This failure is documented in production (Shi et al., 2025) and currently
patched with decoder-side workarounds. We trace the root cause to the retriever and
show that a training-time fix at the encoder eliminates the need for those patches.

Multilingual models such as XLM-RoBERTa (Conneau et al., 2020) learn partially aligned
cross-lingual representations, and contrastive fine-tuning with InfoNCE (van den Oord et al.,
2018) further strengthens this alignment by pulling translation pairs together while
separating non-parallel sentences. Standard retrieval benchmarks such as Tatoeba (Artetxe
and Schwenk, 2019a) and BUCC (Zweigenbaum et al., 2017) evaluate performance by presenting
English queries against a \emph{monolingual} candidate pool---for example, only Spanish
sentences. In this closed-pool setting, semantic alignment and target-language selection
are perfectly correlated by construction: any retrieved sentence is simultaneously the
right meaning and the right language. The two objectives are indistinguishable.

Our experiments expose a gap concealed by this design. After 2000 steps of InfoNCE
fine-tuning on English--Spanish (EN--ES) bitext from OPUS-100, the model achieves
P@1\,=\,0.994 on the standard closed pool of 1012 Spanish candidates. When we expand the
candidate pool to 5060 sentences spanning five languages (ES, FR, DE, SW, AR), all drawn
from FLORES-200 devtest and therefore semantically aligned by construction (each language
contains a translation of the same 1012 source sentences), P@1 drops to 0.382. We
decompose these failures into \emph{wrong-language true positives} (the model retrieves
the semantically correct sentence but in the wrong language) and \emph{true errors}
(the model retrieves an unrelated sentence). The result is unambiguous: P@1-any, where
any correct translation in any language counts as a hit, equals 1.000 for both the encoder
and projected spaces. True errors are exactly zero. Every failure is a wrong-language true
positive.

Language confusion is not uniform and follows a clear typological proximity gradient.
Portuguese, the closest language to Spanish, produces the worst pairwise result
(ES+PT: P@1\,=\,0.364)---worse than the full five-language pool---because all confusion
concentrates on a single competitor. French and German together account for 96\% of
hard-pool failures (ES+FR: 0.567, ES+DE: 0.542), while Swahili and Arabic are near the
easy-pool ceiling (ES+SW: 0.936, ES+AR: 0.902). Notably, German is Germanic, not Romance,
yet confuses nearly as badly as French; the relevant factor is surface-level similarity
(shared script, vocabulary overlap, syntactic proximity), not family membership.

This failure is specific to multiway-parallel evaluation and invisible on standard
benchmarks. When we repeat the hard-pool test on OPUS-100 test data, where each language
pair draws from independent source documents and wrong-language hits are genuine retrieval
errors, the gap drops from 0.612 to 0.013. The model performs normally when candidate
languages do not share source documents. The failure emerges only when semantically
identical sentences in competing languages are present as candidates---precisely the
multiway-parallel setting.

Post-hoc geometric interventions make the hard pool \emph{worse}, not better: whitening
costs 0.039 P@1 points, margin-based scoring (Artetxe and Schwenk, 2019b) costs 0.020,
and both together cost 0.060. Hubness statistics confirm there is no geometry pathology
to fix (hub fraction\,=\,0.000; $N_k$ standard deviation\,=\,1.26 against an expected
value of 2.0). The InfoNCE projection head compounds the problem: the gap in the 256-dim
projected space (0.667) exceeds the gap in the 768-dim encoder space (0.612), meaning the
training objective actively degrades language discrimination rather than merely ignoring
it. Finally, more InfoNCE training worsens target-language selectivity monotonically:
P@1-target drops from 0.382 at 2K steps to 0.259 at 10K before partially recovering to
0.297 at 20K, ruling out undertraining as an explanation and showing that the semantic
collapse InfoNCE induces is progressive.

\textbf{Proposed Research}

We propose a diagnostic and interventional study of target-language selection failure in
bilingual contrastive fine-tuning. Our central thesis is that bilingual InfoNCE optimizes
semantic alignment but leaves target-language specificity underdetermined in multiway-parallel
candidate pools; that this failure is invisible to standard closed-pool evaluation,
follows a typological proximity gradient, and is addressable with a training-time
language-conditioning signal.

First, we \textit{characterize the failure across training pairs}. Beyond the EN--ES pilot,
we fine-tune on EN--FR, EN--DE, and EN--SW and evaluate each against the five-language
hard pool. We test whether confusion severity scales with typological proximity between
the trained language and its neighbors, and whether training on a typologically isolated
pair (EN--SW) produces less cross-lingual confusion than a crowded one (EN--ES).

Second, we \textit{evaluate two training-time interventions}. (i)~\textit{Language-aware
hard negative mining}: we inject same-source wrong-language translations directly into the
InfoNCE denominator using FLORES-200 dev pairs, so the model explicitly learns during
EN--ES training that semantically equivalent French and German sentences are wrong answers.
This requires no architectural change. Hard negative mining closes 99.7\% of the FLORES
gap in 2000 steps. However, the bystander evaluation reveals a severe limitation: by
globally pushing all non-target-language representations away from EN queries during
training, the model collapses EN$\to$FR, EN$\to$DE, and EN$\to$SW easy-pool P@1 to
under 2\%---a catastrophic regression for any system that serves more than one target
language. Standard bilingual EN$\to$ES retrieval also regresses by approximately 10\%
on OPUS-100 test. Hard negative mining is therefore viable only for fixed-target-language
systems, not as a general multilingual retrieval fix. (ii)~\textit{Language-conditioned query encoder}:
motivated directly by the bystander collapse above, we add a learned per-target-language
embedding to the query's mean-pooled representation before the projection head, combined
with an auxiliary contrastive loss over six-tuples (EN, ES, FR, DE, SW, AR) from
FLORES-200 dev that pulls queries toward their target-language translations and pushes
away wrong-language translations of the same sentence. At inference, the query is
conditioned on the desired target language---a targeted pull rather than a global push,
which is the mechanism required to avoid bystander degradation. We sweep auxiliary loss
weight $\lambda \in \{0.1,\,0.5,\,1.0,\,2.0\}$. Because the failure is at least
as severe in the raw encoder space as in the projected space, language conditioning is
applied at the encoder output level to affect the full projection path.

We compare both interventions against: pretrained XLM-R (zero-shot), LaBSE (Feng et al.,
2022) as a zero-shot baseline trained to avoid this collapse, vanilla InfoNCE at the 2K-step
best checkpoint, oracle language filter (ceiling), and predicted language filter
(langdetect). We report P@1-target, P@1-any, and the per-language confusion matrix for
all conditions, on both the training pair (EN--ES) and bystander languages.

Third, we \textit{validate the retrieval fix end-to-end in a RAG pipeline}. Using MKQA
or XQuAD with a multilingual Wikipedia passage corpus (ES, FR, DE, SW, AR), we
demonstrate that the training-time interventions above directly reduce language drift in
generated outputs---confirming that fixing the encoder eliminates the need for decoder-side
workarounds. We measure retrieval language accuracy (fraction of top-1 passages in the
target language) and answer language drift (fraction of LLM outputs in the wrong language).

\textbf{Related Work}

XLM-RoBERTa (Conneau et al., 2020) is our base model. InfoNCE (van den Oord et al., 2018)
and SimCSE (Gao et al., 2021) form the fine-tuning backbone. LaBSE (Feng et al., 2022)
achieves strong cross-lingual retrieval results but trains and evaluates on closed
monolingual pools---the design that conceals the failure mode we study. Artetxe and Schwenk
(2019b) introduce the margin scoring function we test as a post-hoc baseline, and their
LASER system (2019a) establishes the Tatoeba benchmark on which we build.

Shi et al.\ (2025) document language drift in production multilingual RAG
systems---generated output appearing in the wrong language---and propose decoder-side
workarounds; we instead identify and fix the root cause at the retriever.

Huang et al.\ (2025) show that joint cross-lingual alignment and language-specialization
training cause interference and that separate modular training resolves it---directly
validating our failure mode from a complementary angle and motivating our language-conditioned
architecture. Asai et al.\ (2022) demonstrate that same-meaning different-language
translation pairs function as effective hard negatives for multilingual dense retrieval,
supporting our Phase~2a approach.

The closest related work is LAReQA (Roy et al., 2020), which introduces language-agnostic
answer retrieval from a multilingual pool and explicitly notes that closed-pool benchmarks
do not test language confusion bias---directly corroborating our diagnosis. However,
LAReQA optimizes the \emph{opposite} objective: it rewards retrieving the best answer
\emph{regardless} of language. Our setting requires language-targeted retrieval from a
multilingual pool, where semantically correct translations in the wrong language are
failures, not successes. To our knowledge no prior work identifies the gap between
P@1-any and P@1-target as a distinct failure mode of bilingual contrastive training,
characterizes its typological proximity gradient, demonstrates that it vanishes on
non-parallel benchmarks, or proposes a training objective that addresses it without
collapsing retrieval in bystander languages. The bystander collapse produced by hard
negative mining---a known technique---makes the language-conditioned query encoder the
central architectural contribution of this work.

\textbf{Evaluation}

We evaluate retrieval using P@1-target, P@1-any, wrong-language positive rate, true error
rate, and the full per-language confusion matrix under both closed (1012 same-language
candidates) and hard (5060 five-language candidates) pool settings, using FLORES-200
devtest throughout. We additionally report pairwise two-language pool results (ES+PT,
ES+FR, ES+DE, ES+SW, ES+AR) to quantify the proximity gradient and verify that
interventions reduce confusion for the most confusable pairs without degrading the
already-strong SW and AR results. Critically, we evaluate all interventions on
\textit{bystander easy pools}---EN$\to$FR, EN$\to$DE, EN$\to$SW, EN$\to$AR retrieval
on single-language candidate pools---to characterize whether each fix degrades general
multilingual retrieval capability. Hard negative mining collapses bystander P@1 to under
2\%; the language-conditioned encoder is specifically designed to avoid this by
conditioning queries on the target language at inference time rather than globally
repelling bystander languages during training.

We compare FLORES-200 devtest results against OPUS-100 test results to demonstrate the
benchmark design effect: the gap invisible on OPUS-100 (0.013) grows to 0.612 on
FLORES, establishing that multiway-parallel evaluation is necessary to observe the failure
and that the finding is not a generic retrieval failure.

For geometric analysis we track layer-wise CKA and Procrustes similarity before and after
fine-tuning to verify that language conditioning does not degrade the underlying cross-lingual
alignment that standard benchmarks measure. All retrieval experiments are run on the
training pair (EN--ES) and on bystander languages (FR, DE, SW, AR) to confirm that
interventions do not transfer confusion to previously unaffected language pairs.

\textbf{Feasibility and Implementation}

We build on a working codebase using XLM-R base, mean-pooled L2-normalized sentence
representations, and InfoNCE loss with AdamW on OPUS-100 bitext. The hard-pool evaluation
pipeline, error decomposition, per-family pool breakdown, OPUS-100 comparison, and language
confusion matrix are fully implemented and produced the results reported above. Checkpoints
at 2K, 5K, 10K, and 20K steps are trained and available. Multiway parallel data for the
auxiliary loss is available from FLORES-200 dev (997 sentences, separate from the devtest
split used for all reported metrics). The language-conditioned architecture adds one
learned embedding per target language to the existing encoder output; the hard negative
variant modifies only the data collator. Both interventions are feasible on a single GPU
in mixed precision and ablate cleanly against the vanilla InfoNCE baseline. MKQA and
XQuAD are publicly available; the RAG end-to-end evaluation requires one inference pass
per encoder configuration using a publicly available LLM (GPT-4o mini).

\medskip
\noindent\textbf{References}

\noindent Artetxe and Schwenk (2019a). Massively Multilingual Sentence Embeddings for
Zero-Shot Cross-Lingual Transfer and Beyond. \textit{TACL}.

\noindent Artetxe and Schwenk (2019b). Margin-Based Parallel Corpus Mining with
Multilingual Sentence Embeddings. \textit{ACL}.

\noindent Asai et al.\ (2022). Multilingual Dense Passage Retrieval with Hard and False
Negatives. \textit{EMNLP}.

\noindent Conneau et al.\ (2020). Unsupervised Cross-Lingual Representation Learning at
Scale. \textit{EMNLP}.

\noindent Feng et al.\ (2022). Language-Agnostic BERT Sentence Embedding. \textit{ACL}.

\noindent Gao et al.\ (2021). SimCSE: Simple Contrastive Learning of Sentence
Embeddings. \textit{EMNLP}.

\noindent Huang et al.\ (2025). Modular Sentence Encoders: Separating Language
Specialization from Cross-Lingual Alignment. \textit{ACL}. arXiv:2407.14878.

\noindent Roy et al.\ (2020). LAReQA: Language-Agnostic Answer Retrieval from a
Multilingual Pool. \textit{EMNLP}.

\noindent van den Oord et al.\ (2018). Representation Learning with Contrastive Predictive
Coding. \textit{arXiv:1807.03748}.

\noindent Zhang et al.\ (2020). OPUS-100: An English-centric Multilingual Corpus.
\textit{LREC}.

\noindent Shi et al.\ (2025). Language Drift in Multilingual Retrieval-Augmented
Generation. arXiv:2511.09984.

\noindent Zweigenbaum et al.\ (2017). Overview of the Second BUCC Shared Task: Spotting
Parallel Sentences in Comparable Corpora. \textit{ACL Workshop on BUCC}.

\end{document}
